\section{A two-leaf comparison on the whole output interval}\label{cs:two-leaf}

In a completed pure reference the selected leaf range is
\begin{equation}\label{cs:leaf-window}
 \frac14<a<\frac{\arctan2}{\pi}<\frac12.
\end{equation}
A pointwise comparison with a single complete leaf is insufficient
throughout this range. The relevant quantity would be
\[
 \min_{0\le u\le1}\{g(a)+g'(a)(u-a)-k(u)\},
\]
which need not be positive there. The gate joining two leaves supplies the missing comparison. We retain it, and its output tangent, before propagating an estimate into the rest of the tree.

At the selected pure two-leaf stack, write
\begin{equation}\label{cs:seed-parameters}
 X=\cot\pi a,\qquad \cot\pi b=X-\frac1{2X},\qquad
 F_2(b)=2g(a)+m(a,a,b),\qquad
 \lambda=2\log(2X^2).
\end{equation}
Here $\arccot$ takes values in $(0,\pi)$. The displayed expressions define the scalar parameters for every $X\in[1/2,1]$. Completed pure references use the interior range; the three-leaf reference in Section~\ref{cs:probe} uses the endpoint $X=1/2$.

\begin{lemma}[Two-leaf gap]\label{cs:two-leaf-gap}
For $1/2\le X\le1$ and the slope in \eqref{cs:seed-parameters},
\begin{align}
 \max_{D,\,y\le1/2}
 \{k(x)+\widetilde g(y)+m(x,y,z)-\lambda z\}
 &\le F_2(b)-\lambda b-\frac1{100},\label{cs:AC}\\
 \max_D\{k(x)+k(y)+m(x,y,z)-\lambda z\}
 &\le F_2(b)-\lambda b-\frac1{100}.\label{cs:AA}
\end{align}
\end{lemma}

\begin{proof}
Let $P(\lambda)$ be the maximum of $g(x)+g(y)+m(x,y,z)-\lambda z$ on $D$ with $x,y\le1/2$. Direct differentiation at $(a,a,b)$ gives
\[
 m_x=m_y=-2\log X,\qquad m_z=\lambda,\qquad g'(a)=2\log X.
\]
This point is strict in $D$, including the endpoints $X=1/2,1$. Concavity proves
\[
 P(\lambda)=F_2(b)-\lambda b.
\]
The same maximum is obtained if either or both copies of $g$ are replaced by $\widetilde g$: the selected point has $a\ge1/4$, where their supporting tangents agree.

We give the complete finite certificate reduction. Write $H=k+\widetilde g+m$ for the first alias action and $H=k+k+m$ for the second. A rational feasible pure point $p=(p_1,p_1,p_3)$ gives
\[
 P(\lambda)\ge 2g(p_1)+m(p)-\lambda p_3.
\]
For a strict rational alias point $v$, with $v_2>1/4$ in the first case, global concavity gives
$H(w)\le H(v)+\nabla H(v)\cdot(w-v)$. The full vertex sets are
\begin{align*}
 \mathcal V_{\mathrm{AC}}={}&
 \{(0,0,0),(1,0,1),(1/2,1/2,1/2),(0,1/2,1/2),(1/2,1/2,1)\},\\
 \mathcal V_{\mathrm{AA}}={}&
 \{(0,0,0),(1,0,1),(0,1,1),(1/2,1/2,1/2)\}.
\end{align*}
To verify that these lists are complete, the nonnegative coordinates
$p=z-x$, $q_1=z-y$, $r=x+y-z$, $w_0=1-x-y$ obey
$p+q_1+2r+w_0=1$. Thus $D$ is a tetrahedron. Cutting by $y\le1/2$ removes one vertex and adds the two listed edge intersections.

Partition $\lambda/\log2\in[-2,2]$ into the sixteen intervals
$[-2+j/4,-2+(j+1)/4]$, $j=0,\ldots,15$. On each interval the rational data in \cite{Certificates} give one pure point and one tangent point for each alias action. For every vertex $w$ and the two slope endpoints $t$, they give
\begin{equation}\label{cs:two-certificate}
 2g(p_1)+m(p)-t\log2\,p_3
 -H(v)-\nabla H(v)\cdot(w-v)+t\log2\,w_3>\frac1{100}.
\end{equation}
The expression is affine separately in the slope and in $w$. Its endpoint and vertex checks therefore prove the inequality on each full interval and polytope. There are $16\cdot2\cdot(5+4)=288$ inequalities. Section~\ref{cs:arithmetic} proves the outward enclosure rules and specifies the exact data. The smallest certified lower bound is
\begin{equation}\label{cs:two-margin}
 \frac{210002102902309638093785766005979344329}
 {5000000000000000000000000000000000000000}
 >\frac1{25}>\frac1{100}.
\end{equation}
This proves both comparisons, including boundary maxima and the core plateau.
\end{proof}

\section{A new leaf and an arbitrary completed child}\label{cs:attachment}

A completed child is represented by its actual selected output and its output derivative. Its entire internal source remains in the calculation. Let $f_E(b)$ be its optimized pure complete action at output $b$, near the selected strict state. The imported strict maximum has negative Hessian in its internal variables, so the implicit function theorem gives this differentiable local profile. Its derivative is the boundary Lagrange multiplier.

The boundary variables from \cite[Section 6]{Newton} are
\[
 B=\cot\pi b,\qquad
 X^2=\frac{1+B^2}{1+e^{-f_E'(b)}},\qquad
 Y=e^{-f_E'(b)/2}X.
\]
Thus $X^2+Y^2-B^2=1$ and $f_E'(b)=2\log(X/Y)$. At a stack the two children have the same $X$, and their output variables satisfy
\[
 B_o=\frac{B_1B_2+X^2-1}{B_1+B_2},
 \qquad Y_o=\frac{Y_1Y_2}{B_1+B_2}.
\]
For a leaf attached to a completed child these become
\begin{equation}\label{cs:attachment-parameters}
 \begin{gathered}
 a=\arccot X/\pi,\quad b=\arccot B/\pi,\quad
 c=\arccot B_p/\pi,\quad B_p=X-\frac1{X+B},\\
 Y_p=\frac{Y}{X+B},\qquad
 \mu=f_E'(b)=2\log(X/Y),\qquad
 \lambda=f_{hE}'(c)=2\log(X/Y_p).
 \end{gathered}
\end{equation}
For a completed pure reference, both child and parent are completed. Lemma 6.1 of \cite{Newton} then gives $0<X,Y,Y_p<1$. The child inequality $Y<1$ gives $|B|<X$. The parent inequality $Y_p<1$ gives
\[
 B>\frac1{2X}-X.
\]
Together they force $X>1/2$. These two distinct inequalities are needed below.

Define the pure and alias dual maxima
\begin{align}
 P(\mu,\lambda)&=\max_{D,\,x\le1/2}
 \{g(x)+m(x,y,z)+\mu y-\lambda z\},\label{cs:single-P}\\
 Q(\mu,\lambda)&=\max_D
 \{k(x)+m(x,y,z)+\mu y-\lambda z\}.\label{cs:single-Q}
\end{align}
Only the new leaf has the restriction $x\le1/2$. The child coordinate $y$ retains the whole interval $[0,1]$; it is represented by its tangent, not replaced by a second leaf.

\begin{lemma}[Single-attachment gap]\label{cs:attachment-gap}
Let $1/2<X<1$, $Y=\sqrt{1+B^2-X^2}$, $Y_p=Y/(X+B)$ and $0<Y,Y_p<1$, with the positive branch $X+B>0$. Define $a,b,c,\mu,\lambda$ by \eqref{cs:attachment-parameters}. Then
\[
 Q(\mu,\lambda)\le P(\mu,\lambda)-\frac1{2500}.
\]
The pure maximum is attained at $(a,b,c)$ from \eqref{cs:attachment-parameters}. Replacing $g$ by $\widetilde g$ does not change this maximum.
\end{lemma}

\begin{proof}
At the displayed point,
$m_x=-2\log X$, $m_y=-\mu$, $m_z=\lambda$ and $g'(a)=2\log X$. Hence the pure objective is stationary and concavity proves its global maximality. Since $a>1/4$, the same tangent applies to $\widetilde g$.

First suppose $1/2<X\le2/3$. The child gives
$Y^2=1+B^2-X^2\ge1-X^2>X^2$, while $Y<1$. Thus
$-2\log2<2\log X<\mu\le0$. The parent gives separately
$Y_p^2=1+B_p^2-X^2\ge1-X^2>X^2$, while $Y_p<1$, so the same bounds hold for $\lambda$. Both slopes lie in $[-2\log2,0]$.

On the entire closed square, a second rational certificate proves
\begin{equation}\label{cs:square-gap}
 P(\mu,\lambda)-Q(\mu,\lambda)>\frac1{1000}.
\end{equation}
Divide both slopes by $\log2$. The resulting square is partitioned into 55 dyadic rectangles. Each rectangle carries a strict rational pure point $p$, with $p_1<1/2$, and a strict rational tangent point $v$ for $H=k+m$. At each of its four slope corners and each of the four vertices of $D$, the verified inequality is
\begin{multline}\label{cs:single-certificate}
 g(p_1)+m(p)-H(v)-\nabla H(v)\cdot(w-v)\\
 +\mu(p_2-w_2)-\lambda(p_3-w_3)>\frac1{1000}.
\end{multline}
Concavity bounds the alias function by this tangent; the pure point bounds its maximum below. Affinity in the slopes and in $w$ gives the full rectangle and full $D$. The data define each nonterminal subdivision node to have exactly four children, its four quarter-rectangles, starting at $[-2,0]^2$. This proves coverage without gaps or overlap concerns. There are $55\cdot4\cdot4=880$ inequalities. Their least certified lower bound is
\begin{equation}\label{cs:single-margin}
 \frac{55437191760019874892517972684020763373}{10^{40}}
 >\frac1{1000}.
\end{equation}

For $X\ge2/3$ an elementary one-leaf inequality suffices. Set
\[
 I(a)=\min_{0\le u\le1}\{g(a)+g'(a)(u-a)-k(u)\}.
\]
For $a\ge a_*:=\arctan\sqrt2/\pi$, the minimum is at $u=1$ because $g'(a)\le k'(1)=-\log2$. Then
\begin{equation}\label{cs:Ti}
 I(a)=g(a)+(1-a)g'(a)
 =\frac2\pi\Ti(X)+\log X,\qquad
 \Ti(X)=\int_0^X\frac{\arctan t}{t}\dd t.
\end{equation}
For example the two expressions have the same derivative in $X$, namely
$X^{-1}+2\arctan X/(\pi X)$, and their difference tends to zero as $X\downarrow0$. This proves the constant in \eqref{cs:Ti}. The expression decreases with $a$.

For $1/4\le a\le a_*$ the minimizing index is
$u=(2/\pi)\arcsin(\tan^2(\pi a)/2)$. It exceeds $a$: equivalently
$\tan^2 t>2\sin(t/2)$ for $\pi/4\le t\le\arctan\sqrt2$, which holds at $\pi/4$ and remains true because the derivative of the difference is $2\tan t\sec^2t-\cos(t/2)>0$ there and thereafter. Hence $I'(a)=g''(a)(u-a)<0$. The minimum on $1/4\le a\le\arccot(2/3)/\pi$ is consequently at its upper endpoint, where \eqref{cs:Ti} applies.

Four alternating terms of $\Ti(2/3)$ give
\[
 \Ti(2/3)>\frac{1708762}{2679075}.
\]
The series $2\sum_{j\ge0}(1/5)^{2j+1}/(2j+1)$ for $\log(3/2)$, with its geometric tail, gives
$\log(3/2)\le177391/437500$. Together with $\pi<22/7$ these imply
\begin{equation}\label{cs:elementary-margin}
 I(a)>\frac{57269}{136687500}>\frac1{2500}.
\end{equation}
The bound $\pi<22/7$ follows, for instance, from the Machin enclosure used in Section~\ref{cs:arithmetic}. Thus
$k(u)\le g(a)+g'(a)(u-a)-1/2500$ on the full alias interval. Adding the supporting plane of $m$ and the child linear term cancels the input coefficients by stationarity. This proves the lemma in the remaining case.
\end{proof}

\begin{figure}[t]
\centering\begin{tikzpicture}[x=1cm,y=1cm,>=Stealth,
  leaf/.style={draw,rounded corners=2pt,fill=blue!6,minimum width=1.15cm,minimum height=.55cm},
  gate/.style={draw,fill=black!8,minimum size=.38cm,inner sep=1pt},
  every node/.style={font=\small}]
\begin{scope}
 \node[leaf] (h1) at (-1,1.35) {$h$};
 \node[leaf] (h2) at (1,1.35) {$h$};
 \node[gate] (g) at (0,0) {$\star$};
 \draw[->] (h1) -- node[left,pos=.48] {$i$} (g);
 \draw[->] (h2) -- node[right,pos=.48] {$j$} (g);
 \draw[->] (g) -- (0,-1) node[below] {$r$};
 \node[align=center] at (0,-1.8) {$g(a)+g(a)+m(a,a,b)$\\$\lambda=2\log(2X^2)$};
 \node at (0,2) {two-leaf seed};
\end{scope}
\begin{scope}[xshift=6.5cm]
 \node[leaf] (h) at (-1.25,1.35) {$h$};
 \node[draw,rounded corners=2pt,fill=green!5,minimum width=1.8cm,minimum height=1.15cm] (E) at (1.1,1.5) {};
 \node at (1.1,1.9) {$E$};
 \fill (.6,1.55) circle (1pt);
 \fill (1.1,1.55) circle (1pt);
 \fill (1.6,1.55) circle (1pt);
 \draw (.6,1.55)--(.85,1.23)--(1.1,1.55);
 \draw (.85,1.23)--(1.1,.99)--(1.6,1.55);
 \node[gate] (g) at (0,0) {$\star$};
 \draw[->] (h) -- node[left,pos=.48] {$i$} (g);
 \draw[->] (E.south) -- node[right,pos=.48] {$j$} (g);
 \draw[->] (g) -- (0,-1) node[below] {$r$};
 \node[align=center] at (0,-1.8) {$g(a)+f_E(b)+m(a,b,c)$\\$\mu_{\!E}=f_E'(b),\quad\lambda=f_{hE}'(c)$};
 \node at (0,2.55) {single attachment};
\end{scope}
\end{tikzpicture}

\caption{A completed construction retains its exact indices at every step; $\star$ is the multiplication in \eqref{cs:product}. In the left seed, $i,j$ are leaf indices and $r$ is the output of their original multiplication. In the right attachment, $j$ is the output of the entire child $E$, not a new leaf. The affine estimates use the same $i/N,j/N,r/N$; the internal slopes cancel at the selected stationary state.}
\label{cs:figure}
\end{figure}
